\section{Literature Landscape}
\label{sec:related_work}

We organize prior work by \emph{mechanism}, not chronology.
The goal is to show which lines of research our residual-scored RL+GA meta-search depends on (\textbf{used}), and which high-visibility lines we inspected but do not treat as methodological dependencies (\textbf{screened}).
Paraphrases below prefer attached abstracts and OA PDF snippets over title-only name-dropping.

\subsection{Classical discontinuity control and modular audio DSP}
Virtual-analog oscillator and filter design supplies the classical vocabulary for wrap and discontinuity artifacts~\cite{valimaki2006va}.
Aliasing reduction for clipped and discontinuous waveforms, and BLAMP-style corner rounding, motivate seam-local corrections that respect band-limited intent~\cite{esqueda2016aliasing,esqueda2016blamp}.
Differentiable DSP (DDSP) reframes such modules as composable, interpretable blocks inside learning pipelines~\cite{engel2020}; we use that framing for \emph{small} bake/FIR/MLP seam cells rather than end-to-end neural vocoders.
These works are \textbf{used} as domain anchors for the operator vocabulary and for why prolonged cyclic playback is the right outer judge.

\subsection{Label-free and unsupervised restoration}
Noise2Noise shows that restoration mappings can be learned from corrupted observations alone, without clean targets, by exploiting statistical structure of the corruption process~\cite{lehtinen2018}.
That philosophy motivates scoring denoise candidates without paired ``studio-clean'' wavetable cycles: our residual compares an engine-baked tiling to a procedural ideal that withholds open-wrap cliffs, not to a human-cleaned master.
We \textbf{use} Noise2Noise as conceptual justification for unsupervised residual ranking; we do \emph{not} claim to retrain a Noise2Noise CNN on images.

\subsection{Waveform denoising and separation architectures}
Deep learning for audio surveys the move from hand-designed features to learned time/frequency models~\cite{purwins2019}.
Wave-U-Net performs end-to-end source separation with multi-scale 1-D U-Nets~\cite{stoller2018}; singing-voice U-Nets similarly popularized encoder--decoder skips on audio~\cite{jansson2017}.
Conv-TasNet demonstrated strong time-domain separation with temporal convolutional architectures~\cite{luo2019}.
These lines are \textbf{used} as design priors for \emph{tiny} searchable cells (shallow U-Net / dilated / dual-path / attention blocks under a strict per-trial bake budget).
They are not used as frozen full-scale separators inside every meta trial.

\subsection{Neural architecture search and RL controllers}
Elsken et al.\ taxonomize NAS by search space, strategy, and evaluation cost~\cite{elsken2019}.
ENAS couples an RL controller with parameter sharing for efficient discrete search~\cite{pham2018}.
We \textbf{use} this family as the conceptual ancestor of RL proposing architecture edits (add/remove/mutate ops, toggle residual-primary, mutate MLP/FIR), while scoring with prolonged residual rather than ImageNet reward.
PPO supplies the clipped-surrogate policy update used when the overnight loop logs PPO mutation branches~\cite{schulman2017}.

\subsection{Evolutionary NAS and population schedules}
Regularized / aging evolution showed that evolutionary NAS can discover competitive image classifiers when population turnover is carefully regularized~\cite{real2019}.
Population Based Training jointly adapts weights and hyperparameters in an asynchronous population, discovering schedules rather than a single fixed hyperparameter vector~\cite{jaderberg2017}.
Practical Bayesian optimization~\cite{snoek2012} and decomposition-based multi-objective evolution (MOEA/D)~\cite{zhang2007} further inform local Bayesian and multi-objective prior families in earlier DenoiseOpt searches.
These are \textbf{used} for GA tournament/crossover branches, PBT-style exploit--mutate, and literature-combo priors---implemented at synthesizer-trial scale, not as distributed datacenter PBT.

\subsection{Evolutionary reinforcement learning hybrids}
Khadka and Tumer's evolution-guided policy gradient (ERL) interleaves an evolutionary population with an RL actor to mitigate sparse reward, weak exploration, and brittle hyperparameters in deep RL~\cite{khadka2018}.
We \textbf{use} ERL as the \emph{spirit} of interleaving GA population steps with PPO updates at tiny population sizes; we do not claim a full ERL robotics stack.

\subsection{Mixture-of-experts soft gating}
Sparsely gated MoE layers increase capacity via conditional computation~\cite{shazeer2017}.
We \textbf{use} MoE as motivation for soft gates over heterogeneous seam blocks (e.g., U-Net vs dilated vs attention cells) under residual scoring---again at bake-scale widths, not trillion-parameter routing.

\subsection{Used versus screened}
\label{sec:used-screened}

\paragraph{Used (methodological or comparative anchors).}
\begin{itemize}
  \item Discontinuity / VA lineage:~\cite{valimaki2006va,esqueda2016aliasing,esqueda2016blamp}.
  \item Unsupervised restoration philosophy:~\cite{lehtinen2018}.
  \item Modular / differentiable audio DSP context:~\cite{engel2020}.
  \item Audio DL and compact waveform cells:~\cite{purwins2019,stoller2018,jansson2017,luo2019}.
  \item NAS / RL controllers / PPO:~\cite{elsken2019,pham2018,schulman2017}.
  \item Evolutionary NAS and population HPO:~\cite{real2019,jaderberg2017,snoek2012,zhang2007}.
  \item ERL hybrid outer loop:~\cite{khadka2018}.
  \item MoE soft gating prior:~\cite{shazeer2017}.
\end{itemize}

\paragraph{Screened (relevant contrast; not dependencies).}
\begin{itemize}
  \item \textbf{MAML}~\cite{finn2017}: inspirational bi-level / fast-adaptation framing; we do \emph{not} differentiate through a deep task network. Inner fits are coordinate-style updates on bake parameters; outer rank is residual $R$.
  \item \textbf{DARTS}~\cite{liu2019}: continuous relaxation of architecture search; we keep discrete ops + small cells under residual scoring.
  \item \textbf{Demucs / Hybrid Demucs}~\cite{dfossez2019,dfossez2021}: strong music separators, but full multi-scale / hybrid STFT stacks are too heavy per meta trial; screened for overnight bake budgets.
  \item \textbf{DiffWave}~\cite{kong2021}: powerful diffusion vocoder; screened as a generative sampling stack, not a seam bake operator.
  \item \textbf{SEGAN-style adversarial enhancement}~\cite{pascual2017}: optional tiny adversarial cues only; full GAN training loops screened out for per-trial cost.
\end{itemize}

\paragraph{Positioning.}
Relative to this landscape, the present work contributes (i)~a prolonged residual objective for periodic seam denoising, (ii)~an explicitly hybrid RL+GA(+PBT) meta-outer loop with honest branch naming, and (iii)~a used/screened literature protocol tied to OA PDF analyses.
We do not claim SOTA on general speech enhancement or music separation corpora; the primary claim is protocol-level: residual-scored meta-search for cycle-local audio denoise operators under a declared compute budget.
